\chapter{Blepharitis}
\index{Blepharitis}
\label{sec:Blepharitis}

\section{Overview}
Blepharitis is a chronic inflammatory condition affecting the eyelid margins. It is classified into anterior blepharitis (involving the base of the eyelashes) and posterior blepharitis (involving the meibomian glands). It is a common, often recurrent disorder that can cause significant ocular discomfort and may contribute to dry eye disease.

\section{Classification}
\begin{description}[font=\normalfont\bfseries,leftmargin=3em,labelwidth=2.5em]
  \item[Cause] Bacterial colonization (anterior) or meibomian gland dysfunction (posterior)
  \item[Organ System] Eyes
  \item[Pathophysiology] Anterior blepharitis: staphylococcal infection or seborrheic inflammation of lash follicles. Posterior blepharitis: meibomian gland obstruction with altered lipid secretion, inflammatory mediator release, and tear film instability.
  \item[Duration] Chronic
  \item[Onset] Gradual
  \item[Mode of Transmission] Not directly transmissible
  \item[Clinical Course] Chronic relapsing-remitting course with periods of exacerbation and quiescence.
  \item[Severity] Mild to moderate; rarely severe
  \item[Extent of Disease] Localized to eyelid margins
  \item[Age Group] More common in older adults; can occur at any age
  \item[Epidemiology] Affects 37--47\% of ophthalmology patients; prevalence increases with age. Associated with rosacea, seborrheic dermatitis, and dry eye syndrome.
  \item[ICD-11 Code] 9A02.0
\end{description}

\section{Causes}
Anterior blepharitis is caused by staphylococcal infection or seborrheic dermatitis at the eyelash follicles. Posterior blepharitis results from meibomian gland dysfunction with abnormal lipid secretion, leading to inflammation and tear film instability.

\section{Risk Factors}
\begin{itemize}[noitemsep]
  \item Seborrheic dermatitis or rosacea
  \item Dry eye syndrome
  \item Contact lens wear
  \item Eyelash mites (\emph{Demodex folliculorum})
  \item Diabetes mellitus
  \item Immunosuppression
\end{itemize}

\section{Signs and Symptoms}
\begin{itemize}[noitemsep]
  \item Itching, burning, or gritty sensation in the eyes
  \item Erythema and swelling of eyelid margins
  \item Crusting and collarettes at the base of eyelashes
  \item Tearing or dry eye symptoms
  \item Sensitivity to light (photophobia)
  \item Difficulty performing daily activities
\end{itemize}

\section{Progression and Stages}
Blepharitis follows a chronic relapsing-remitting pattern. Exacerbations may last days to weeks and are triggered by stress, poor hygiene, or environmental factors. Without consistent management, the condition may lead to meibomian gland atrophy, eyelash loss (madarosis), or corneal complications.

\section{Complications}
\begin{itemize}[noitemsep]
  \item Dry eye syndrome (evaporative dry eye)
  \item Recurrent hordeola or chalazia
  \item Conjunctivitis or keratitis
  \item Corneal ulceration or scarring (rare)
  \item Eyelash loss or trichiasis
  \item Reduced quality of life
  \item Emotional and psychological effects
\end{itemize}

\section{Diagnosis}
\begin{itemize}[noitemsep]
  \item Slit-lamp examination of eyelid margins and lashes
  \item Meibomian gland expression and assessment
  \item Lab tests (microbial culture in refractory cases)
  \item Imaging tests (not routinely needed)
\end{itemize}

\section{Prevention}
\begin{itemize}[noitemsep]
  \item Daily eyelid hygiene with warm compresses and lid scrubs
  \item Avoid eye rubbing and sharing cosmetics
  \item Manage underlying skin conditions (rosacea, seborrhea)
  \item Regular check-ups
\end{itemize}

\section{Treatment Options}
\begin{itemize}[noitemsep]
  \item Warm compresses followed by eyelid margin massage
  \item Eyelid hygiene with diluted baby shampoo or commercial lid wipes
  \item Topical antibiotics (erythromycin, bacitracin) or combination antibiotic-steroid ointment
  \item Oral doxycycline or azithromycin for posterior blepharitis
  \item Artificial tears for associated dry eye
  \item Lifestyle changes
\end{itemize}

\section{Living with It}
\begin{itemize}[noitemsep]
  \item Follow treatment plan
  \item Regular appointments
  \item Adjust daily habits
  \item Support groups
  \item Keep learning
\end{itemize}

\section{Outlook}
Blepharitis is a chronic condition that cannot be cured but can be effectively controlled with consistent eyelid hygiene and appropriate medical therapy. Long-term management is usually required to prevent exacerbations and complications. Most patients maintain good visual function with proper care.
